\documentclass{report}
\usepackage{amsmath,amssymb}
\setlength{\parindent}{0mm}
\setlength{\parskip}{1em}
\begin{document}
\begin{center}
\rule{6in}{1pt} \
{\large James Nagy \\
{\bf Separable Nonlinear Inverse Problems and the Local Minimum Trap}}

Mathematics and Computer Science \\ Emory University \\ Atlanta \\ GA 30322
\\
{\tt nagy@mathcs.emory.edu}\\
Anastasia Cornelio\\
Elena Loli Piccolomini\end{center}

In this talk we consider a nonlinear least squares framework to solve
separable nonlinear ill-posed inverse problems. It is shown that with
proper constraints and well chosen regularization parameters, it is
possible to obtain an objective function that is fairly well behaved.
Although uncertainties in the data and inaccuracies of linear solvers
make it unlikely to obtain a smooth and convex objective function, it
is shown that implicit filtering optimization methods can be used to
avoid becoming trapped in local minima. An application to blind
deconvolution is used for illustration.


\end{document}
